% id: 116-19
% book: 베이직쎈 확률과 통계 · 06 이항분포와 정규분포
% section: 유형 072 정규분포의 활용; 도수 구하기
% figure: tikz:fig-116-19
% answer: 
% answer_source: 
% variant_level: 0 (원본 전사)
% 이 파일은 build.mjs 가 items.json 에서 만든다. 고칠 때는 items.json 을 고친다.
\smash{\raisebox{1.5mm}{\dmkichul{베이직쎈 확률과 통계 116쪽 19번}}}
\noindent\begin{minipage}[t]{0.58\linewidth}\vspace{0pt}\raggedright 어느 양계장에서 생산되는 달걀 100개의 무게는 평균이 $52\mkern3mu \mathrm{g}$, 표준편차가 $3\mkern3mu \mathrm{g}$인 정규분포를 따른다고 한다. 오른쪽 표준정규분포표를 이용하여 무게가 $46\mkern3mu \mathrm{g}$ 이상 $55\mkern3mu \mathrm{g}$ 이하인 달걀의 개수를 구하면?\end{minipage}\hfill\begin{minipage}[t]{0.4\linewidth}\vspace{0pt}\centering \resizebox{\ifdim\width>\linewidth\linewidth\else\width\fi}{!}{% 116-19(인쇄 116쪽) — 116-19 오른쪽에 놓인 표준정규분포표($z$ / $\mathrm{P}(0\le Z\le z)$).
% 스캔 표라 크롭하지 않고 tabular 로 다시 조판했다(칸 값은 원문 그대로 · 원문에 없는 행은 넣지 않는다).
{\setlength{\tabcolsep}{5pt}\renewcommand{\arraystretch}{1.25}%
\begin{tabular}{|c|c|}
\hline
$z$ & $\mathrm{P}(0\le Z\le z)$ \\
\hline
$1.0$ & $0.34$ \\
\hline
$1.5$ & $0.43$ \\
\hline
$2.0$ & $0.48$ \\
\hline
\end{tabular}}
}\end{minipage}\par
\dmchoicesauto{}{$78$}{$80$}{$82$}{$84$}{$86$}
